% !TITLE 大风歌
% !AUTHOR 刘邦
% !DYNASTY 两汉
% !GENRE 楚歌
% !ORDER 1
% !EDITION 《史记》，司马迁撰，中华书局（1959年点校本）
% !EDITION_STATUS pending

\begin{poem}
大风起兮\textsuperscript{1}云飞扬\textsuperscript{2}，
威加海内\textsuperscript{3}兮归故乡，
安得猛士\textsuperscript{4}兮守四方\textsuperscript{5}！
\end{poem}

\begin{annotations}
\notetext{1}{兮：语气助词，楚辞和楚歌中常用的停顿，相当于今天口语里的“啊”。}
\notetext{2}{云飞扬：云气翻卷飞扬，写大风骤起时云层被卷动的景象。}
\notetext{3}{威加海内：威名施加于四海之内，即威震天下。海内，四海之内，指全国。}
\notetext{4}{安得猛士：怎样才能得到勇猛的将士。安，怎么、怎样；得，得到。}
\notetext{5}{四方：东西南北，指天下四方，这里指帝国的边境。}
\end{annotations}

\begin{appreciation}
刘邦传世的诗作只有两首，都是楚歌。一首《鸿鹄歌》唱给戚夫人，一首《大风歌》唱给故乡。汉高祖十二年，他平定英布叛乱，回师途中经过沛县，在沛宫置酒，召来故人父老纵情饮宴。酒到酣处，他亲自击筑，唱出了这首歌。

开头“大风起兮云飞扬”——大风骤起，云气翻卷。这是酒席上抬眼可见的实景，也是他半生征战的写照：秦末天下大乱，群雄并起，风云际会之中，一个沛县亭长一路走到了帝国的最高处。“兮”是楚辞的句法，刘邦是楚人，唱的是楚地的调子。屈原在《离骚》里用“兮”字写忧愤，刘邦用“兮”字写豪情，同一个语气词，隔着一百多年，装下了两种截然不同的人生。

第二句“威加海内兮归故乡”——威名加于四海之内，然后回到故乡。“威加海内”把整个天下收进一句，“归故乡”又把天下收回到一个具体的地方。衣锦还乡是中国人最古老的荣耀，刘邦唱得理所当然，没有半点遮掩。

可就在这份得意之后，第三句忽然转冷：“安得猛士兮守四方！”怎样才能得到勇猛的将士，替我守住四方？天下已经打下来了，守天下才刚刚开始。当年替他打天下的韩信、彭越、英布，一个个或死或叛；北边的匈奴虎视眈眈，而异姓诸侯王已被剪除殆尽，能替他守四方的猛士，反而找不到了。唱完这首歌，他起身起舞，据《史记》记载，“慷慨伤怀，泣数行下”。

三句之间有一个巨大的落差。前两句越唱越高，第三句忽然落回地面。帝王的威风背后，是一个开国者无人可托的孤独。

《大风歌》最动人的地方，是威风之后的那一问。得到天下的人，在故乡的酒席上，想到的却是失去。几个月后，刘邦去世。这首三句的歌，成了他留给故乡最后的声响。
\end{appreciation}
